% paper1_discussion.tex -- Discussion and Conclusion of Paper 1, input by
% papers/paper1_indicators.tex.
\section{Discussion}

The index ranked the three systems in the same order at every site, which is
the behaviour a soil health index needs if it is to guide management: the
integrated system, which returns the most organic matter to the soil, scored
highest, and the conventional system lowest. The ranking was driven mainly by
the biological and physical indicators. Soil organic carbon, aggregate
stability and microbial biomass all respond to organic inputs and reduced
disturbance, and their scores moved together, as the links between soil
organisms, organic matter and structure would lead one to expect
\parencite{doran2000}.

The index also separated the sites more strongly than the systems. The
semi-arid site scored lowest under every system, because its low carbon and
high bulk density kept several indicators near the bottom of their scoring
range. This illustrates a known limitation of indices built on fixed scoring
limits: they mix the inherent capacity of a soil with the effect of
management, so a sandy soil in a dry climate can appear unhealthy even when it
is well managed \parencite{bunemann2018}. Comparing systems within a site, as
done here, avoids that confusion; comparing sites requires scoring limits that
reflect the potential of each soil, as the Soil Management Assessment
Framework proposes \parencite{andrews2004}.

Two further cautions apply. Equal weighting assumes that every indicator is
equally important to the functions of interest, which is a choice rather than
a finding. And soil pH, scored on an optimum range, contributed little to the
differences between systems because it changed little within the duration of
the trial; its role in the index is to flag soils that move out of the range
in which nutrients are available.

\section{Conclusion}

A minimum dataset of six indicators, converted to scores and averaged into a
single index, ranked three management systems consistently at three
contrasting sites, with the index ranging from \res{p1-shi-min} to
\res{p1-shi-max}. The index is best used to compare management within a site;
comparison across sites requires scoring limits matched to each soil. All
values in this paper come from synthetic data and illustrate the method only.
